\subsection{忠实平坦环}

命题 8. 设 A、B 是两个环， \(\rho\) 是 A 到 B 的同态。假设存在右 B-模 E 使 \(\rho_{*}(E)\) 是忠实平坦 A-模。则：

\begin{enumerate}
\def\labelenumi{(\roman{enumi})}
\item
  对每个左 A-模 F，典范同态 \(j: \mathbf{F} \to \mathbf{F}_{(\mathbf{B})} = \mathbf{B} \otimes_{\mathbf{A}} \mathbf{F}\) （它满足 \(j(x) = 1 \otimes x\) ，对所有 \(x \in \mathbf{F}\) ）是单射。
\item
  对 A 的每个左理想 \(\mathfrak{a}\) ，有 \(\mathbf{A},\overline{\rho}^{-1}(\mathbf{Ba}) = \mathfrak{a}\) 。
\item
  同态 \(\rho\) 是单射。
\item
  对 A 的每个极大左理想 \(\mathfrak{m}\) ，存在 B 的极大左理想 \(\mathfrak{n}\) 使 \(\overline{\rho}^{-1}(\mathfrak{n}) = \mathfrak{m}\) 。
\end{enumerate}

我们证明 (i)。我们知道（《代数学》第二章 §5 第 2 目命题 5 的推论），对每个右 B-模 M，典范 A-同态 \(i: \mathbf{M} \to \rho_*(\mathbf{M}) \otimes_{\mathbf{A}} \mathbf{B} = \rho^*(\rho_*(\mathbf{M}))\) （由 \(i(y) = y \otimes 1\) 定义）是单射，且 A-模 \(i(\mathbf{M})\) 是 \(\rho_*(\mathbf{M}) \otimes_{\mathbf{A}} \mathbf{B}\) 的直因子。因此，对每个左 A-模 F，

\[
i \otimes 1 _ {F}: \rho_ {*} (M) \otimes_ {A} F \rightarrow \rho_ {*} (M) \otimes_ {A} B \otimes_ {A} F
\]

是单射（§2 第 1 目，引理 2）。取 \(\mathbf{M} = \mathbf{E}\) ，则（由于 \(i \otimes 1_{\mathbf{F}} = 1_{\mathbf{M}} \otimes j\) ）推出 \(j\) 是单射（第 1 目，命题 2）。

断言 (ii) 由 (i) 取 \(\mathbf{F} = \mathbf{A}_s / \mathfrak{a}\) 得出，而 (iii) 由 (ii) 取 \(\mathfrak{a} = \{0\}\) 得出。

最后，若 \(\mathfrak{m}\) 是 \(\mathbf{A}\) 的极大左理想，则由 (ii) 有 \(\bar{\rho}^{-1}(\mathbf{B}\mathfrak{m}) = \mathfrak{m}\) ，因此 \(\mathbf{B}\mathfrak{m} \neq \mathbf{B}\) 。于是存在 \(\mathfrak{n}\) ，即 \(\mathbf{B}\) 的一个包含 \(\mathbf{B}\mathfrak{m}\) 的极大左理想（《代数学》第一章 §8 第 7 目，定理 2）；这时 \(\mathfrak{m} \subset \bar{\rho}^{-1}(\mathfrak{n})\) ，且由于 \(\rho(1) \notin \mathfrak{n}\) ，有 \(1 \notin \bar{\rho}^{-1}(\mathfrak{n})\) 。因此 \(\bar{\rho}^{-1}(\mathfrak{n}) = \mathfrak{m}\) 。

若 A 与 B 满足命题 8 的条件，通常借助 \(\rho\) 把 A 等同于 B 的一个子环。

推论. 在命题 8 的假设下，若 B 是左 Noether 的（相应地左 Artin 的），则 A 亦然。

若 \((\mathfrak{a}_n)\) 是 A 的左理想组成的非稳定递增（相应地递减）序列，则 B 的理想序列 \((\mathrm{Ba}_n)\) 是非稳定且递增（相应地递减）的，因为 \(\overline{\rho}^{-1}(\mathrm{Ba}_n) = \mathfrak{a}_n\) ，这与假设相矛盾。

*注 (1). 若 A 与 B 交换，我们将在第二章 §2 第 5 目命题 11 的推论 4 中看到，命题 8 的假设蕴含

对每个 \(\text{prime}\) 理想 \(\mathfrak{p}\) （属于 \(\mathbf{A}\) ），存在 \(\text{prime}\) 理想 \(\mathfrak{q}\) （属于 \(\mathbf{B}\) ）使 \(\overline{\rho}^{-1}(\mathfrak{q}) = \mathfrak{p}\) （其中 \(\mathfrak{p} = \mathbf{A} \cap \mathfrak{q}\) ，当 \(\mathbf{A}\) 等同于 \(\mathbf{B}\) 的一个子环时）。*

命题 8 的一个重要应用是 B 本身为忠实平坦 A-模的情形。但在这种情形，我们有下面更精确的命题：

命题 9. 设 A、B 是两个环， \(\rho\) 是 A 到 B 的同态。下列五个性质等价：

\begin{enumerate}
\def\labelenumi{(\alph{enumi})}
\item
  右 A-模 B 是忠实平坦的。
\item
  同态 \(\rho\) 是单射，且右 A-模 \(\mathbf{B} / \rho (\mathbf{A})\) 是平坦的。
\item
  右 A-模 B 是平坦的，且对每个左 A-模 F，典范同态 \(x \mapsto 1 \otimes x\) （从 F 到 B \(\otimes_{A} F\) ）是单射。
\item
  右 A-模 B 是平坦的，且对 A 的每个左理想 a， \(\overline{\rho}^{-1}(\mathbf{B}\mathfrak{a}) = \mathfrak{a}\) 。
\item
  右 A-模 B 是平坦的，且对 A 的每个极大左理想 m，存在
\end{enumerate}

B 的一个极大左理想 n 使 \(\bar{\rho}^{-1}(\mathfrak{n}) = \mathfrak{m}\) 。

由命题 8，(a) 蕴含性质 (c)、(d)、(e) 中的每一个。另一方面，若 (e) 成立，则 \(\mathbf{Bm} \neq \mathbf{B}\) 对 A 的每个极大左理想 \(m\) 成立（因为存在 B 的极大左理想 \(n\) 使 \(\mathbf{Bm} \subset n\) ），且由第 1 目命题 1 的判别准则 (d)，B 是忠实平坦 A-模；因此 (e) 蕴含 (a)。

我们现在证明 \((c) \Rightarrow (d) \Rightarrow (b) \Rightarrow (a)\) ，从而完成证明。首先，在 (c) 中取 \(F = A_{s}/a\) 便知 (c) 蕴含 (d)。若 (d) 成立，取 \(a = \{0\}\) 便得 \(\rho\) 是单射；(d) 与 §2 第 6 目命题 7 的推论蕴含 \(B/\rho(A)\) 是平坦右 A-模，即 (d) 蕴含 (b)。最后，若 (b) 成立，把第 1 目的命题 3 应用于正合序列

\[
0 \longrightarrow \mathrm{A} _ {d} \stackrel {\rho} {\longrightarrow} \mathrm{B} \longrightarrow \mathrm{B} / \rho (\mathrm{A}) \longrightarrow 0
\]

由于 \(A_{d}\) 是忠实平坦的，这表明 B 是忠实平坦右 A-模。

*注 (2). 若 A 与 B 交换，我们将在第二章 §2 第 5 目命题 11 的推论 4 中看到，命题 9 的诸条件等价于下列条件：

\begin{enumerate}
\def\labelenumi{(\alph{enumi})}
\setcounter{enumi}{5}
\tightlist
\item
  B 是平坦 A-模，且对 A 的每个素理想 p，存在 B 的理想 q 使 \(\overline{\rho}^{-1}(\mathfrak{q}) = \mathfrak{p}\) 。
\end{enumerate}

在命题 9 的条件下，借助 \(\rho\) 把 A 等同于 B 的一个子环。这时关系式 \(\bar{\rho}^{-1}(\mathbf{Ba}) = a\) 就读作 \(\mathbf{A} \cap \mathbf{Ba} = a\) 。另一方面，若 F 是左 A-模，则 F 等同于它在 \(\mathbf{B} \otimes_{\mathbf{A}} \mathbf{F}\) 中的在典范映射 \(x \mapsto 1 \otimes x\) 下的像；若 X 是 F 的加法子群，我们用 BX 表示由 X 生成的 \(\mathbf{B} \otimes_{\mathbf{A}} \mathbf{F}\) 的左 B-子模。采用这个记号，我们有：

命题 10. 设 B 是环，A 是 B 的子环，且 B 是忠实平坦右 A-模。设 F 是左 A-模，F'、F'' 是 F 的两个子模。则：

\begin{enumerate}
\def\labelenumi{(\roman{enumi})}
\item
  典范映射 \(\mathbf{B} \otimes_{\mathbf{A}} \mathbf{F}' \to \mathbf{B} \otimes_{\mathbf{A}} \mathbf{F}\) 是 \(\mathbf{B} \otimes_{\mathbf{A}} \mathbf{F}'\) 到 \(\mathbf{BF}'\) 上的同构。
\item
  \(\mathbf{F} \cap \mathbf{BF}' = \mathbf{F}'\) 。
\item
  \(\mathrm{B}(\mathbf{F}' + \mathbf{F}'') = \mathbf{BF}' + \mathbf{BF}''\) 。
\item
  \(\mathbf{B}(\mathbf{F}'\cap \mathbf{F}'')=\mathbf{B}\mathbf{F}'\cap \mathbf{B}\mathbf{F}''\) 。
\end{enumerate}

由于 B 是平坦右 A-模，典范映射

\[
\mathbf {B} \otimes_ {\mathbf {A}} \mathbf {F} ^ {\prime} \rightarrow \mathbf {B} \otimes_ {\mathbf {A}} \mathbf {F}
\]

是单射；考虑到所作的等同，它的像是 BF'，这就证明了 (i)。断言 (ii) 由 §2 第 6 目命题 7 取 E = B、E' = A 并利用公式 \(A \otimes_{A} F = F\) 与 \(A \otimes_{A} F' = F'\) 得出。断言 (iii) 是平凡的，而 (iv) 由 §2 第 6 目命题 6 得出。

